\chapter{The Gamma and Beta Functions}\label{ch:gamma}

\begin{center}\itshape
$1!=1$, $2!=2$, $3!=6$, $4!=24$.\\
Which curve passes through $\Gamma(n+1)=n!$, and why is there only one?
\end{center}

On 13 October 1729 Euler wrote to Goldbach with a product that returns $n!$ at the positive integers and still makes sense for a fraction. On 8 January 1730 he sent an integral for the same job. Daniel Bernoulli had written to Goldbach a week earlier, on 6 October 1729, with a product of his own. The integral is the one this chapter studies.

Many smooth curves pass through the points $(n,(n-1)!)$. Bohr and Mollerup, in 1922, proved that only one of them has a convex logarithm. That uniqueness is the reason $\Gamma$ is the interpolation. Artin's book of 1931 turns the theorem into three checks: the functional equation, a convex logarithm, and the value $1$ at $1$.

The checks give Gauss's product, the Beta integral, and Legendre's duplication formula. With the sine product of Chapter~\ref{ch:series} they give Euler's reflection formula $\Gamma(x)\Gamma(1-x)=\frac\pi{\sin\pi x}$. At the end, the volume of the unit ball in $\R^n$ is $\frac{\pi^{n/2}}{\Gamma(\frac n2+1)}$. For $n=5$ that volume is $\frac{8\pi^2}{15}\approx5.2638$. It tends to $0$ as $n$ grows.

\begin{think}
Suppose the factorial had a natural value at $\frac12$. What should $\bigl(\frac12\bigr)!$ be? A first guess is $\bigl(\frac12\bigr)!^2\approx0!\cdot1!$. Proposition~\ref{prop:gammahalf} gives the value.
\end{think}

\section{Euler's integral}

The integral below is the curve. At the positive integers it has to return the factorials, one step down.

\begin{definition}
For $x>0$, the \emph{Gamma function} is
\[
\Gamma(x)=\int_0^\infty t^{x-1}e^{-t}\dd t.
\]
\end{definition}

The integral converges. Near $0$ the integrand is at most $t^{x-1}$, and that is integrable because $x>0$. Near infinity it is $O(t^{-2})$, because $t^{x+1}e^{-t}\to0$.

\begin{proposition}\label{prop:gammafe}
$\Gamma(1)=1$ and $\Gamma(x+1)=x\Gamma(x)$ for $x>0$. Hence $\Gamma(n+1)=n!$ for every integer $n\ge0$.
\end{proposition}
\begin{proof}
$\Gamma(1)=\int_0^\infty e^{-t}\dd t=1$. Integrate by parts on $[\varepsilon,R]$:
\[
\int_\varepsilon^Rt^xe^{-t}\dd t=\Bigl[-t^xe^{-t}\Bigr]_\varepsilon^R+x\int_\varepsilon^Rt^{x-1}e^{-t}\dd t,
\]
The boundary term tends to $0$ as $\varepsilon\to0$ and $R\to\infty$. Induction gives the values at the positive integers.
\end{proof}

So $\Gamma(2)=1!$, $\Gamma(3)=2!$, $\Gamma(4)=3!$ and $\Gamma(5)=4!$. The values are $1$, $2$, $6$ and $24$. Figure~\ref{fig:rw05-factorials} draws the curve through those four points. The last one stands far above the first three.

\begin{figure}[ht]
\centering
\fig{ch05-rewrite-factorials}
\caption{$\Gamma$ passes through $1!$, $2!$, $3!$ and $4!$ at $x=2,3,4,5$. The values are $1$, $2$, $6$ and $24$.}
\label{fig:rw05-factorials}
\end{figure}

\begin{proposition}\label{prop:gammasmooth}
$\Gamma$ is of class $C^\infty$ on $\left]0,\infty\right[$, with
\[
\Gamma^{(k)}(x)=\int_0^\infty t^{x-1}(\ln t)^ke^{-t}\dd t.
\]
\end{proposition}
\begin{proof}
Fix $0<a<A$. For $x\in[a,A]$ and $t>0$, the bound $t^{x-1}\le t^{a-1}+t^{A-1}$ holds: the first term is the larger one for $t\le1$, and the second for $t\ge1$. Hence
\[
\bigl|t^{x-1}(\ln t)^ke^{-t}\bigr|\le\bigl(t^{a-1}+t^{A-1}\bigr)|\ln t|^ke^{-t},
\]
The right-hand side is continuous and integrable on $\left]0,\infty\right[$ for every $k$. Theorem~\ref{thm:domination}, applied $k$ times on $J=[a,A]$, says $\Gamma$ is $C^k$ there, with that derivative. The numbers $a$, $A$ and $k$ were arbitrary, so the claim holds on $\left]0,\infty\right[$.
\end{proof}

The functional equation pushes $\Gamma$ to the left of $0$. For $x\in\left]-n,0\right[$ not an integer, set
\[
\Gamma(x)=\frac{\Gamma(x+n)}{x(x+1)\cdots(x+n-1)}.
\]
The value does not depend on $n$: replacing $n$ by $n+1$ multiplies the numerator and the denominator by $x+n$. The relation
$\Gamma(x+1)=x\Gamma(x)$ still holds. Near $-n$, with $n\ge0$ an integer,
\[
\Gamma(x)=\frac{\Gamma(x+n+1)}{x(x+1)\cdots(x+n)}\sim\frac{(-1)^n}{n!\,(x+n)},
\]
since the numerator tends to $\Gamma(1)=1$ and $x(x+1)\cdots(x+n-1)\to(-n)(-n+1)\cdots(-1)=(-1)^nn!$. So $\Gamma$ has a
simple pole at each of $0,-1,-2,\dots$ (Figure~\ref{fig:gamma}). It never vanishes. On $\left]0,\infty\right[$ it is the integral of a
positive function, and the formula above carries that fact to the left. Chapter~\ref{ch:complex} extends $\Gamma$ to the complex plane.

\begin{figure}[ht]
\centering
\fig{ch05-gamma}
\caption{The real graph of $\Gamma$. The curve passes through the points $(n,(n-1)!)$. On $\left]0,\infty\right[$ the minimum is at $x_0\approx1.4616$, where $\Gamma(x_0)\approx0.8856$. There is a pole at every integer $n\le0$, and the function changes sign there.}
\label{fig:gamma}
\end{figure}

\section{Log-convexity and the Bohr--Mollerup theorem}

A positive function $f$ is \emph{log-convex} when $\ln f$ is convex. Log-convexity is a stronger demand than convexity of $f$, because $f=e^{\ln f}$ is then a convex increasing function of a convex function, and $e^u$ is convex and increasing. Products, positive powers and translations $x\mapsto x+c$ preserve it, since $\ln(fg)=\ln f+\ln g$.

The next lemma is why $\Gamma$ is log-convex. The inequality is first an inequality of non-negative functions, so it survives the limit from compact pieces of the integral.

\begin{lemma}\label{lem:logconvex}
Let $I$ be an interval, $u\ge0$ and $\ell$ continuous on $I$, and $J$ an interval such that $F(x)=\int_Iu(t)e^{x\ell(t)}\dd t$
converges and is positive for every $x\in J$. Then $F$ is log-convex on $J$.
\end{lemma}
\begin{proof}
\emph{One inequality at each $t$.}
Let $x,y\in J$, $\lambda\in\left]0,1\right[$ and $p=\lambda x+(1-\lambda)y$. For $A,B\ge0$, the weighted AM--GM inequality (the concavity
of $\ln$, Volume~1, Chapter~9) gives $A^\lambda B^{1-\lambda}\le\lambda A+(1-\lambda)B$. Apply it, for each $t$, to
\[
\begin{gathered}A=\frac{u(t)e^{x\ell(t)}}{F(x)},\qq B=\frac{u(t)e^{y\ell(t)}}{F(y)},\end{gathered}
\]
for which $A^\lambda B^{1-\lambda}=\frac{u(t)e^{p\ell(t)}}{F(x)^\lambda F(y)^{1-\lambda}}$. Both sides are non-negative. On each compact subinterval of $I$ the inequality integrates. The integrals that define $F(x)$, $F(y)$ and $F(p)$ converge, so the same comparison survives the limit over an exhaustion of $I$. Integrating over $I$ therefore gives
$\frac{F(p)}{F(x)^\lambda F(y)^{1-\lambda}}\le\lambda+(1-\lambda)=1$, that is, $\ln F(p)\le\lambda\ln F(x)+(1-\lambda)\ln F(y)$.
\end{proof}

\begin{corollary}\label{cor:gammalogconvex}
$\Gamma$ is log-convex on $\left]0,\infty\right[$.
\end{corollary}
\begin{proof}
Write $\Gamma(x)=\int_0^\infty\frac{e^{-t}}t\,e^{x\ln t}\dd t$ and apply Lemma~\ref{lem:logconvex} with $u(t)=\frac{e^{-t}}t$ and $\ell(t)=\ln t$.
\end{proof}

A convex function $\varphi$ has one property we use over and over. The slope $s(a,b)=\frac{\varphi(b)-\varphi(a)}{b-a}$ of a chord is a non-decreasing
function of each endpoint. Indeed, for $a<b<c$, write $b=\lambda a+(1-\lambda)c$ with $\lambda=\frac{c-b}{c-a}$; the convexity inequality
$\varphi(b)\le\lambda\varphi(a)+(1-\lambda)\varphi(c)$ rearranges both to $s(a,b)\le s(a,c)$ and to $s(a,c)\le s(b,c)$. In particular, for
$a<b\le c<d$, $s(a,b)\le s(b,c)\le s(c,d)$ (with the middle term omitted if $b=c$).

\begin{theorem}[Bohr--Mollerup]\label{thm:bohrmollerup}
Let $f:\left]0,\infty\right[\to\left]0,\infty\right[$ satisfy $f(1)=1$, $f(x+1)=xf(x)$ for all $x>0$, and be log-convex. Then $f=\Gamma$.
Moreover, for every $x>0$, Gauss's product formula holds:
\[
\Gamma(x)=\lim_{n\to\infty}\frac{n!\,n^x}{x(x+1)\cdots(x+n)}.
\]
\end{theorem}
\begin{proof}
\emph{The chords.}
Let $\varphi=\ln f$, which is convex, with $\varphi(x+1)=\varphi(x)+\ln x$; by induction $\varphi(n+1)=\ln n!$ and
\[
\varphi(x+n+1)=\varphi(x)+\ln\bigl(x(x+1)\cdots(x+n)\bigr).
\]
Let $0<x\le1$ and $n\ge1$. The three intervals $[n,n+1]$, $[n+1,n+1+x]$ and $[n+1,n+2]$ have chords of increasing slope
(the first lies to the left of the second, and the second is the third shortened on the right), so
\[
\ln n\le\frac{\varphi(n+1+x)-\varphi(n+1)}x\le\ln(n+1).
\]
\emph{The gap.}
Multiply by $x>0$. The two formulas for $\varphi$ turn the middle expression into $\varphi(x)+\ln\bigl(x(x+1)\cdots(x+n)\bigr)-\ln n!$. So
\[
x\ln n\le\varphi(x)+\ln\bigl(x(x+1)\cdots(x+n)\bigr)-\ln n!\le x\ln(n+1).
\]
Set
\[
g_n(x)=\ln\frac{n!\,n^x}{x(x+1)\cdots(x+n)}.
\]
The left inequality is $g_n(x)\le\varphi(x)$. The right inequality is $\varphi(x)\le g_n(x)+x\ln\bigl(1+\frac1n\bigr)$, because $\ln(n+1)-\ln n=\ln\bigl(1+\frac1n\bigr)$. The result is $0\le\varphi(x)-g_n(x)\le x\ln\bigl(1+\frac1n\bigr)$. Hence $\varphi(x)=\lim_ng_n(x)$ for $0<x\le1$: the hypotheses determine $f$
on $\left]0,1\right]$, and then on $\left]0,\infty\right[$ through $f(x+1)=xf(x)$. Since $\Gamma$ satisfies the hypotheses
(Proposition~\ref{prop:gammafe}, Corollary~\ref{cor:gammalogconvex}), $f=\Gamma$.

\emph{Every positive $x$.}
Gauss's formula is proved for $0<x\le1$. A direct computation gives
\[
g_n(x+1)-g_n(x)=\ln x+\ln\frac n{x+n+1},
\]
which tends to $\ln x$. So if $g_n(x)\to\ln\Gamma(x)$, then $g_n(x+1)\to\ln\Gamma(x)+\ln x=\ln\Gamma(x+1)$, and the formula holds for all $x>0$ by induction.
\end{proof}

For a numerical check, take $x=1$ in Gauss's product. The fraction is $\frac{n!\,n}{1\cdot2\cdots(1+n)}=\frac{n}{n+1}$, which tends to $1=\Gamma(1)$.

\begin{think}
The function $f(x)=\Gamma(x)\bigl(1+\frac3{10}\sin2\pi x\bigr)$ is positive, with $f(1)=1$ and $f(x+1)=xf(x)$. Which hypothesis of Theorem~\ref{thm:bohrmollerup} fails, and on which interval? See Figure~\ref{fig:loggamma}.
\end{think}

\begin{figure}[ht]
\centering
\fig{ch05-loggamma}
\caption{Charcoal: $\ln\Gamma$, which is convex. Coral: $\ln\Gamma(x)+\ln\bigl(1+\frac3{10}\sin2\pi x\bigr)$. The coral curve agrees with $\ln\Gamma$ at the integers and obeys the same functional equation. The wiggles destroy convexity.}
\label{fig:loggamma}
\end{figure}

The wiggle fails log-convexity, and that is the hypothesis which picks $\Gamma$ out of the other curves through the factorials.

\begin{corollary}[Weierstrass's product]\label{cor:weierstrass}
For every $x>0$, with $\gamma$ Euler's constant (Volume~1, Chapter~3),
\[
\frac1{\Gamma(x)}=xe^{\gamma x}\prod_{k\ge1}\Bigl(1+\frac xk\Bigr)e^{-x/k},
\]
and the product converges.
\end{corollary}
\begin{proof}
Let $H_n=\sum_{k\le n}\frac1k$ and rewrite $n^{-x}=e^{x(H_n-\ln n)}e^{-xH_n}$. Then
\[
\begin{aligned}\frac{x(x+1)\cdots(x+n)}{n!\,n^x}&=xe^{x(H_n-\ln n)}\eqbr\cdot\prod_{k=1}^n\Bigl(1+\frac xk\Bigr)e^{-x/k}.\end{aligned}
\]
Gauss's formula and the limit $H_n-\ln n\to\gamma$ send the left side to $\frac1{\Gamma(x)}$. The partial products therefore converge to $\frac{e^{-\gamma x}}{x\Gamma(x)}\ne0$.
\end{proof}

At $x=1$ the product says $\frac1{\Gamma(1)}=e^{\gamma}\prod_{k\ge1}\bigl(1+\frac1k\bigr)e^{-1/k}$. The partial product up to $n$ is $e^{\gamma}\cdot\frac{n+1}1\cdot e^{-H_n}$, which tends to $1$ because $H_n-\ln(n+1)\to\gamma$. That is the same limit, written the other way around.

\section{Stirling's formula for the Gamma function}

Volume~1 proved $n!\sim\sqrt{2\pi n}\bigl(\frac ne\bigr)^n$ (Chapter~8; the shape with an unspecified constant is in Chapter~3). Log-convexity fills the gaps between the integers.

\begin{theorem}\label{thm:stirlinggamma}
As $x\to\infty$ through real values, $\Gamma(x)\sim\sqrt{2\pi}\,x^{x-\frac12}e^{-x}$.
\end{theorem}
\begin{proof}
\emph{Between two integers.}
Write $x=n+t$ with $n$ an integer and $t\in[0,1[$, and let $\varphi=\ln\Gamma$. Convexity at the point $n+t=(1-t)n+t(n+1)$ gives
$\varphi(n+t)\le\varphi(n)+t\ln n$. At the point $n+1=t(n+t)+(1-t)(n+1+t)$ it gives
$\varphi(n+1)\le t\varphi(n+t)+(1-t)\varphi(n+1+t)$, and $\varphi(n+1+t)=\varphi(n+t)+\ln(n+t)$, so
\[
\varphi(n+1)\le\varphi(n+t)+(1-t)\ln(n+t).
\]
Together, with $\Gamma(n+1)=n\Gamma(n)$,
\[
\Bigl(\frac n{n+t}\Bigr)^{1-t}\le\frac{\Gamma(n+t)}{\Gamma(n)\,n^t}\le1,
\]
and the left side is at least $\frac n{n+1}$. So $\Gamma(x)=\Gamma(n)n^t(1+o(1))$, uniformly in $t$.
\emph{Stirling's formula for $n!$.}
By Stirling's formula for $n!$,
$\Gamma(n)=\frac{n!}n\sim\sqrt{2\pi}\,n^{n-\frac12}e^{-n}$. On the other side,
\[
\begin{aligned}x^{x-\frac12}e^{-x}&=n^{n-\frac12}e^{-n}\,n^t\eqbr\cdot\Bigl(1+\frac tn\Bigr)^n e^{-t}\cdot\Bigl(1+\frac tn\Bigr)^{t-\frac12}.\end{aligned}
\]
\emph{The two expansions.}
The logarithm of the middle factor is $n\ln\bigl(1+\frac tn\bigr)-t$. For $u\ge0$, $\ln(1+u)\le u$, so this logarithm is at most $0$. Also $h(u)=\ln(1+u)-u+\frac{u^2}2$ has $h(0)=0$ and $h'(u)=\frac{u^2}{1+u}\ge0$, so $\ln(1+u)\ge u-\frac{u^2}2$ and the logarithm is at least $-\frac{t^2}{2n}\ge-\frac1{2n}$. The last factor lies between
$(1+\frac1n)^{-1/2}$ and $(1+\frac1n)^{1/2}$. Comparing the two expansions proves the theorem.
\end{proof}

At $n=4$, $\Gamma(4.5)/\Gamma(4)=\frac{35\sqrt\pi}{32}\approx1.93862$, while $\sqrt4=2$. The ratio is already close to the square root that Stirling predicts for a step of size $\frac12$.

Chapter~\ref{ch:asymptotics} continues the expansion as Stirling's series $\ln\Gamma(x)=(x-\frac12)\ln x-x+\frac12\ln2\pi+\frac1{12x}-\frac1{360x^3}+\cdots$. The series is asymptotic, and it diverges at every $x$.

\section{The Beta integral}

Beta is the integral on $[0,1]$ that the same three checks will identify with a ratio of Gamma values.

\begin{definition}
For $x,y>0$, the \emph{Beta function} is
\[
B(x,y)=\int_0^1t^{x-1}(1-t)^{y-1}\dd t.
\]
\end{definition}

The integral converges since $x,y>0$, and $B(x,y)=B(y,x)$ by $t\mapsto1-t$. For positive integers, $B(m,n)$ is the reciprocal of
an integer multiple of a binomial coefficient. The next theorem says why. For a first check, $B(2,3)=\int_0^1 t(1-t)^2\dd t=\frac1{12}$.

\begin{lemma}\label{lem:betarec}
$B(1,y)=\frac1y$ and $B(x+1,y)=\frac x{x+y}B(x,y)$ for $x,y>0$.
\end{lemma}
\begin{proof}
The identity $B(1,y)=\int_0^1(1-t)^{y-1}\dd t=\frac1y$ is the first claim. For the second, integrate by parts on $[\varepsilon,1-\varepsilon]$ and let $\varepsilon\to0$. The boundary terms vanish because $x,y>0$:
\[
\begin{aligned}B(x+1,y)&=\int_0^1t^x(1-t)^{y-1}\dd t\eqbr=\frac xy\int_0^1t^{x-1}(1-t)^y\dd t\eqbr=\frac xyB(x,y+1).\end{aligned}
\]
Expand $(1-t)^y=(1-t)^{y-1}-t(1-t)^{y-1}$. Then $B(x,y+1)=B(x,y)-B(x+1,y)$. Substitute and solve for $B(x+1,y)$.
\end{proof}

\begin{theorem}[Euler]\label{thm:betagamma}
For $x,y>0$,
\[
B(x,y)=\frac{\Gamma(x)\Gamma(y)}{\Gamma(x+y)}.
\]
\end{theorem}
\begin{proof}
\emph{The candidate.}
Fix $y>0$ and let $f(x)=\frac{\Gamma(x+y)B(x,y)}{\Gamma(y)}$ for $x>0$. Then $f>0$, and $f(1)=\frac{\Gamma(1+y)}{y\Gamma(y)}=1$ by
Lemma~\ref{lem:betarec}. With the lemma and $\Gamma(x+1+y)=(x+y)\Gamma(x+y)$,
\[
\begin{aligned}f(x+1)&=\frac{(x+y)\Gamma(x+y)}{\Gamma(y)}\cdot\frac x{x+y}B(x,y)\eqbr=xf(x).\end{aligned}
\]
\emph{Log-convexity, then uniqueness.}
Finally $x\mapsto B(x,y)=\int_0^1\frac{(1-t)^{y-1}}t\,e^{x\ln t}\dd t$ is log-convex by Lemma~\ref{lem:logconvex}, $x\mapsto\Gamma(x+y)$ is a
translate of a log-convex function, and $\Gamma(y)$ is a constant, so $f$ is log-convex. By Theorem~\ref{thm:bohrmollerup}, $f=\Gamma$.
\end{proof}

At $(x,y)=(2,3)$ the formula says $B(2,3)=\frac{\Gamma(2)\Gamma(3)}{\Gamma(5)}=\frac{1\cdot2}{24}=\frac1{12}$, which matches the integral above.

\begin{remark}
A second proof writes $\Gamma(x)\Gamma(y)$ as an integral over the quarter plane. That argument is in Volume~7. The proof above uses only the uniqueness theorem.
\end{remark}

\begin{proposition}\label{prop:gammahalf}
$\Gamma\bigl(\frac12\bigr)=\sqrt\pi$, so $\bigl(\frac12\bigr)!=\Gamma\bigl(\frac32\bigr)=\frac{\sqrt\pi}2\approx0.886$.
\end{proposition}
\begin{proof}
By Theorem~\ref{thm:betagamma}, $\Gamma(\frac12)^2=\Gamma(1)B(\frac12,\frac12)=\int_0^1\frac{\mathrm dt}{\sqrt{t(1-t)}}$. With $t=\sin^2\theta$,
$\mathrm dt=2\sin\theta\cos\theta\dd\theta$ and $\sqrt{t(1-t)}=\sin\theta\cos\theta$ on $[0,\frac\pi2]$, so the quotient is $2$ and the integral is $\int_0^{\pi/2}2\dd\theta=\pi$. Since $\Gamma(\frac12)>0$, $\Gamma(\frac12)=\sqrt\pi$.
\end{proof}

The opening guess was $\sqrt{0!\cdot1!}=1$. The value $\frac{\sqrt\pi}2\approx0.886227$ is smaller, and log-convexity is why the guess had to sit above the curve. The substitution $t=u^2$ turns $\Gamma(\frac12)$ into $2\int_0^\infty e^{-u^2}\dd u$, so this is the Gaussian integral again. It is the third evaluation in the series: Wallis integrals in Volume~1, Chapter~8, the parameter in Example~\ref{ex:gaussfeynman}, and this substitution.

\section{Reflection and duplication}

Reflection ties $\Gamma(x)$ to $\Gamma(1-x)$. The sine product from Chapter~\ref{ch:series} is the one extra ingredient.

\begin{theorem}[Euler's reflection formula]\label{thm:reflection}
For every real $x$ that is not an integer,
\[
\Gamma(x)\Gamma(1-x)=\frac\pi{\sin\pi x}.
\]
\end{theorem}
\begin{proof}
\emph{On the unit interval.}
First let $0<x<1$. Let $G_n(x)=\frac{n!\,n^x}{x(x+1)\cdots(x+n)}$, so that $G_n(x)\to\Gamma(x)$ by Gauss's formula, and the same at $1-x$.
In the product $G_n(x)G_n(1-x)$ the powers of $n$ combine to $n^1$. In the denominator, group the factors $k+x$ and $k-x$ for
$1\le k\le n$, which gives $\prod_k(k^2-x^2)=(n!)^2\prod_k\bigl(1-\frac{x^2}{k^2}\bigr)$; the factors $x$ and $n+1-x$ are left over. So
\[
G_n(x)G_n(1-x)=\frac n{n+1-x}\cdot\frac1{x\prod_{k=1}^n\bigl(1-\frac{x^2}{k^2}\bigr)}.
\]
As $n\to\infty$ the right side tends to $\frac\pi{\sin\pi x}$ by the sine product, Theorem~\ref{thm:sineproduct}.

\emph{The same value on every interval.}
For general $x$, let $R(x)=\Gamma(x)\Gamma(1-x)\sin\pi x$. By the functional equation,
$\Gamma(x+1)\Gamma(-x)=x\Gamma(x)\cdot\frac{\Gamma(1-x)}{-x}$, and $\sin\pi(x+1)=-\sin\pi x$, so $R(x+1)=R(x)$. Since $R=\pi$ on $\left]0,1\right[$,
$R=\pi$ on every interval $\left]n,n+1\right[$.
\end{proof}

At $x=\frac12$ the formula gives $\Gamma(\frac12)^2=\pi$ once more. At $x=\frac16$, $\sin\frac\pi6=\frac12$, so $\Gamma(\frac16)\Gamma(\frac56)=2\pi\approx6.283185$. At $x=\frac14$ it gives $\Gamma(\frac14)\Gamma(\frac34)=\pi\sqrt2$, which
Exercise~\ref{ex:lemniscate} uses to evaluate the arc length of the lemniscate.

\begin{theorem}[Legendre's duplication formula]\label{thm:duplication}
For $x>0$,
\[
\Gamma(x)\Gamma\Bigl(x+\frac12\Bigr)=2^{1-2x}\sqrt\pi\,\Gamma(2x).
\]
\end{theorem}
\begin{proof}
\emph{Three checks.}
Let $f(x)=\frac{2^{x-1}}{\sqrt\pi}\Gamma\bigl(\frac x2\bigr)\Gamma\bigl(\frac{x+1}2\bigr)$. Then $f>0$, and $f(1)=\frac1{\sqrt\pi}\Gamma(\frac12)\Gamma(1)=1$.
Next,
\[
\begin{aligned}f(x+1)&=\frac{2^x}{\sqrt\pi}\Gamma\Bigl(\frac{x+1}2\Bigr)\Gamma\Bigl(\frac x2+1\Bigr)\eqbr=\frac{2^x}{\sqrt\pi}\Gamma\Bigl(\frac{x+1}2\Bigr)\frac x2\Gamma\Bigl(\frac x2\Bigr)\eqbr=xf(x).\end{aligned}
\]
Finally $\ln f$ is a sum of an affine function and two convex functions of the form $x\mapsto\ln\Gamma(ax+b)$ with $a>0$, so
it is convex. By Theorem~\ref{thm:bohrmollerup}, $f=\Gamma$. Replacing $x$ by $2x$ gives the formula.
\end{proof}

At $x=1$ both sides equal $\frac{\sqrt\pi}2$. The left side is $\Gamma(1)\Gamma(\frac32)=\frac{\sqrt\pi}2$. The right side is $2^{-1}\sqrt\pi\,\Gamma(2)=\frac{\sqrt\pi}2$.

Exercise~\ref{ex:multiplication} runs the same argument for Gauss's multiplication formula. There are $m$ factors in place of two.

\section{The volume of the unit ball}

Let $V_n$ be the volume of the unit ball $\{x\in\R^n:\ x_1^2+\cdots+x_n^2\le1\}$, with $V_0=1$. By scaling, the ball of radius $r$
has volume $r^nV_n$. Slicing the $n$-ball perpendicular to the last coordinate, the slice at height $t\in[-1,1]$ is an
$(n-1)$-ball of radius $\sqrt{1-t^2}$, so
\[
V_n=\int_{-1}^1V_{n-1}\bigl(1-t^2\bigr)^{\frac{n-1}2}\dd t.
\]
This is Cavalieri's principle: a volume is the integral of the areas of its slices. Here that formula is the definition of $V_n$. It agrees with the familiar $V_2=\pi$ and $V_3=\frac43\pi$. Volume~7 proves the same formula for Lebesgue measure.

\begin{theorem}\label{thm:ballvolume}
For every $n\ge0$,
\[
V_n=\frac{\pi^{n/2}}{\Gamma\bigl(\frac n2+1\bigr)}.
\]
In particular $V_{2m}=\frac{\pi^m}{m!}$, and $V_n\to0$ as $n\to\infty$.
\end{theorem}
\begin{proof}
\emph{One slice.}
By symmetry and the substitution $t=\sqrt s$, for $a>-1$,
\[
\int_{-1}^1(1-t^2)^a\dd t=2\int_0^1(1-t^2)^a\dd t.
\]
Put $t=\sqrt s$, so $\dd t=\frac{\dd s}{2\sqrt s}$. The factor $2$ cancels the $2$ in the denominator, and
\[
\begin{aligned}\int_{-1}^1(1-t^2)^a\dd t&=\int_0^1s^{-1/2}(1-s)^a\dd s\eqbr=B\Bigl(\frac12,a+1\Bigr)\eqbr=\frac{\Gamma(\frac12)\Gamma(a+1)}{\Gamma(a+\frac32)}.\end{aligned}
\]
\emph{The product of the ratios.}
With $a=\frac{n-1}2$ and $\Gamma(\frac12)=\sqrt\pi$, the slicing formula becomes
\[
\frac{V_n}{V_{n-1}}=\sqrt\pi\,\frac{\Gamma(\frac{n+1}2)}{\Gamma(\frac n2+1)}.
\]
The product of these ratios for $1\le k\le n$ telescopes. The numerator at step $k$ is $\Gamma(\frac{k+1}2)$, and the denominator at step $k-1$ is $\Gamma(\frac{k-1}2+1)=\Gamma(\frac{k+1}2)$. The numerator at $k=1$ is $\Gamma(1)$, and the denominator at $k=n$ is $\Gamma(\frac n2+1)$. Since $V_0=1$,
$V_n=\pi^{n/2}\frac{\Gamma(1)}{\Gamma(\frac n2+1)}$. Finally $\Gamma(\frac n2+1)$ grows faster than any $C^n$, by Theorem~\ref{thm:stirlinggamma}, so
$V_n\to0$.
\end{proof}

For the even case, $n=2m$ gives $\Gamma(m+1)=m!$ and $V_{2m}=\frac{\pi^m}{m!}$. In particular $V_4=\frac{\pi^2}2\approx4.934802$ and $V_6=\frac{\pi^3}6\approx5.167713$. The odd values need half-integers. One finds $V_1=2$, $V_3=\frac{4\pi}3\approx4.188790$ and $V_5=\frac{8\pi^2}{15}\approx5.263789$. Then $V_{10}=\frac{\pi^5}{120}\approx2.550164$ and $V_{20}=\frac{\pi^{10}}{10!}\approx0.025807$. The five-dimensional ball is about $204$ times the twenty-dimensional one: $V_5/V_{20}\approx203.968$.

\begin{figure}[ht]
\centering
\fig{ch05-ball}
\caption{Volumes $V_n$ of the unit ball in $\R^n$. The sequence rises through $V_5=\frac{8\pi^2}{15}\approx5.264$ and then falls toward $0$, with $V_{20}\approx0.026$.}
\label{fig:ball}
\end{figure}

\begin{figure}[ht]
\centering
\fig{ch05-rewrite-balls}
\caption{The unit ball has volumes $2$, $3.1416$, $4.1888$, $5.2638$, $2.5502$ and $0.02581$ in dimensions $1$, $2$, $3$, $5$, $10$ and $20$. The volume peaks at dimension $5$ and then collapses.}
\label{fig:rw05-balls}
\end{figure}

The ratio of successive volumes drops below $1$ between $n=5$ and $n=6$. Indeed $V_5/V_4=\frac{16}{15}$ and $V_6/V_5=\frac{5\pi}{16}$. Since $\pi<\frac{16}5$, the second ratio is less than $1$.

\begin{think}
The unit ball sits inside the cube $[-1,1]^n$, which has volume $2^n$, and the ball meets every face. Even so $\frac{V_n}{2^n}\to0$ very fast. Where does the volume of the cube go? Compare the distance from the centre to a corner with the radius of the ball.
\end{think}

The maximum at $n=5$ depends on the radius being $1$: the ball of radius $r$ has volume $r^nV_n$, and the ratio $\frac{V_n}{V_{n-1}}$ decreases like $\sqrt{2\pi/n}$. The collapse does not depend on that choice. In high dimension almost all of the cube lies in its corners, far from the inscribed ball, and almost all of the ball lies in a thin shell near its surface. The fraction of the volume within distance $\varepsilon$ of the surface is $1-(1-\varepsilon)^n$. For $n=20$ and $\varepsilon=\frac1{10}$, that fraction is $1-(0.9)^{20}\approx0.878$. This is the geometry behind the failure of grid methods and rejection sampling in high dimension.

\section{What we proved, and what comes next}

Bohr and Mollerup say that $\Gamma$ is the only positive function on $]0,\infty[$ with $f(x+1)=xf(x)$, $f(1)=1$ and a convex logarithm (Theorem~\ref{thm:bohrmollerup}). Each new formula was three checks, and uniqueness finished the proof. Gauss's product, the Beta integral $B(x,y)=\frac{\Gamma(x)\Gamma(y)}{\Gamma(x+y)}$ (Theorem~\ref{thm:betagamma}) and Legendre's duplication formula (Theorem~\ref{thm:duplication}) came out that way. The reflection formula (Theorem~\ref{thm:reflection}) needed the sine product of Chapter~\ref{ch:series} as well. Stirling's formula moved from the integers to every large real $x$ (Theorem~\ref{thm:stirlinggamma}). The volume of the unit ball is $\frac{\pi^{n/2}}{\Gamma(\frac n2+1)}$ (Theorem~\ref{thm:ballvolume}), and $V_n\to0$.

So $\bigl(\frac12\bigr)!=\Gamma\bigl(\frac32\bigr)=\frac{\sqrt\pi}2\approx0.886$. The first guess, $\sqrt{0!\cdot1!}=1$, was a little too large, and log-convexity is why.

The closed form for $V_n$ is order. The collapse, with almost all of the cube out in the corners, is the surprise that high dimension forces on a uniform draw.

Gamma returns later: Chapter~\ref{ch:complex} extends it to the complex plane, Chapter~\ref{ch:zeta} puts it in the functional equation of $\zeta$, Chapter~\ref{ch:asymptotics} finds the full Stirling series, and Chapter~\ref{ch:agm} computes $\Gamma(\frac14)$ from the arithmetic--geometric mean. The next chapter asks a different question. Can a sum of sines draw a function that jumps?
\section{Exercises}

\begin{exercise}[$\star$]\label{ex:halfint}
Prove $\Gamma\bigl(n+\frac12\bigr)=\frac{(2n)!}{4^nn!}\sqrt\pi$ for every integer $n\ge0$, and compute $\Gamma\bigl(-\frac12\bigr)$ and $\Gamma\bigl(-\frac32\bigr)$.
\end{exercise}

\begin{exercise}[$\star$]\label{ex:expp}
Prove $\int_0^\infty e^{-x^p}\dd x=\Gamma\bigl(1+\frac1p\bigr)$ for $p>0$. What is the limit as $p\to\infty$, and why is that limit obvious?
\end{exercise}

\begin{exercise}[$\star\star$]\label{ex:sincos}
For $a,b>-1$, prove
\[
\int_0^{\pi/2}\sin^a\theta\cos^b\theta\dd\theta=\frac12B\Bigl(\frac{a+1}2,\frac{b+1}2\Bigr),
\]
and recover the Wallis integrals $W_n=\int_0^{\pi/2}\sin^n\theta\dd\theta$ of Volume~1, Chapter~8.
\end{exercise}

\begin{exercise}[$\star\star$]\label{ex:sphere}
Define the area $S_{n-1}$ of the unit sphere in $\R^n$ as the derivative at $r=1$ of the volume $r^nV_n$ of the ball of radius $r$, the volume of a thin shell divided by its thickness. Prove that $S_{n-1}=nV_n=\frac{2\pi^{n/2}}{\Gamma(n/2)}$, check $S_1=2\pi$ and $S_2=4\pi$, and
find the $n$ for which $S_{n-1}$ is largest.
\end{exercise}

\begin{exercise}[$\star\star$]\label{ex:lemniscate}
The lemniscate $(x^2+y^2)^2=x^2-y^2$ has total length $2\varpi$, where $\varpi=2\int_0^1\frac{\mathrm dr}{\sqrt{1-r^4}}$ (Chapter~\ref{ch:agm}). Prove that
\[
\int_0^1\frac{\mathrm dr}{\sqrt{1-r^4}}=\frac{\Gamma(\frac14)^2}{4\sqrt{2\pi}}.
\]
\end{exercise}

\begin{exercise}[$\star\star\star$]\label{ex:multiplication}
Establish Gauss's multiplication formula, for every integer $m\ge1$ and every $x>0$:
\[
\prod_{k=0}^{m-1}\Gamma\Bigl(x+\frac km\Bigr)=(2\pi)^{\frac{m-1}2}m^{\frac12-mx}\,\Gamma(mx).
\]
\end{exercise}

\begin{exercise}[$\star\star\star$]\label{ex:digamma}
Prove $\Gamma'(1)=-\gamma$, that is $\int_0^\infty e^{-t}\ln t\dd t=-\gamma$, where $\gamma$ is Euler's constant.
\end{exercise}

\begin{exercise}[$\star\star$]\label{pb:nair}
\begin{enumerate}[label=(\alph*)]
\item Let $d_n=\operatorname{lcm}(1,2,\dots,n)$. Prove that $m\binom nm$ divides $d_n$ for $1\le m\le n$.
\item Conclude that $d_n\ge2^n$ for $n\ge7$, and that $\pi(n)\ge\frac{n\log2}{\log n}$ for $n\ge7$.
\end{enumerate}
\end{exercise}

\begin{exercise}[$\star\star\star$]\label{pb:ncycles}
Let $n\ge2$ and let $c=(1\,2\,\dots\,n)$. Let $V\subset\C^n$ be the hyperplane $\{\sum x_i=0\}$. The group $\mathfrak S_n$ permutes the coordinates. Write $\chi_r$ for the character of the exterior power $\Lambda^rV$, $0\le r\le n-1$. Quoted, and not used by any theorem of this volume: the
$\Lambda^rV$ are irreducible and pairwise non-isomorphic (Fulton and Harris, \emph{Representation Theory}, Proposition~3.12).
\begin{enumerate}[label=(\alph*)]
\item Prove that $\chi_r(c)=(-1)^r$, and that every other irreducible character vanishes at $c$.
\item Prove that the number of ways to write $g=ab$ with $a$ in a class $A$ and $b$ in a class $B$ is
$\frac{|A||B|}{|G|}\sum_\chi\frac{\chi(a)\chi(b)\overline{\chi(g)}}{\chi(1)}$, in any finite group.
\item Prove that $\sum_{r=0}^m\frac{(-1)^r}{\binom mr}=\frac{m+1}{m+2}\bigl(1+(-1)^m\bigr)$.
\item Conclude that $c$ is the product of two $n$-cycles in exactly $\frac{2(n-1)!}{n+1}$ ways if $n$ is odd, and in none if $n$ is
even.
\end{enumerate}
\end{exercise}

\begin{exercise}[$\star$]\label{pb:dalzell}
\emph{Why $\pi<\frac{22}7$.}
\begin{enumerate}[label=(\alph*)]
\item Prove that $\displaystyle\int_0^1\frac{x^4(1-x)^4}{1+x^2}\dd x=\frac{22}7-\pi$, and deduce that $\frac1{1260}<\frac{22}7-\pi<\frac1{630}$.
\item Prove that for every $n\ge1$ there is a rational $r_n$ with $\displaystyle\int_0^1\frac{x^{4n}(1-x)^{4n}}{1+x^2}\dd x=r_n+(-4)^n\frac\pi4$, and that
$\bigl|\pi+4(-4)^{-n}r_n\bigr|<4^{1-5n}$.
\end{enumerate}
\end{exercise}

\begin{exercise}[$\star$]\label{pb:sophomore}
\emph{The sophomore's dream (Johann Bernoulli, 1697).}
\begin{enumerate}[label=(\alph*)]
\item Prove that $\displaystyle\int_0^1x^{-x}\dd x=\sum_{n\ge1}\frac1{n^n}$ and $\displaystyle\int_0^1x^x\dd x=\sum_{n\ge1}\frac{(-1)^{n+1}}{n^n}$.
\item How many terms fix the first integral to $10$ decimal places?
\end{enumerate}
\end{exercise}

\begin{exercise}[$\star$]\label{ex:rw05a}
From $\Gamma(1)=1$ and $\Gamma(x+1)=x\Gamma(x)$, compute $\Gamma(6)$. Then compute $B(3,2)$ from Theorem~\ref{thm:betagamma}, and check it by evaluating $\int_0^1 t^2(1-t)\dd t$.
\end{exercise}

\begin{exercise}[$\star\star$]\label{ex:rw05b}
Use Theorem~\ref{thm:reflection} to show that $\Gamma\bigl(\frac16\bigr)\Gamma\bigl(\frac56\bigr)=2\pi$ and $\Gamma\bigl(\frac13\bigr)\Gamma\bigl(\frac23\bigr)=\frac{2\pi}{\sqrt3}$. Give both products as decimals.
\end{exercise}

\begin{exercise}[$\star\star$]\label{ex:rw05c}
From Theorem~\ref{thm:ballvolume}, show that $V_1=2$, $V_2=\pi$, $V_3=\frac{4\pi}3$, $V_4=\frac{\pi^2}2$, $V_5=\frac{8\pi^2}{15}$, $V_6=\frac{\pi^3}6$, $V_{10}=\frac{\pi^5}{120}$ and $V_{20}=\frac{\pi^{10}}{10!}$. Show that $\frac{V_5}{V_4}=\frac{16}{15}$ and $\frac{V_6}{V_5}=\frac{5\pi}{16}$. Since $\pi<\frac{16}5$, deduce $V_6<V_5$. Compute the six decimals of $V_1,V_2,V_3,V_5,V_{10},V_{20}$, and the ratio $V_5/V_{20}$.
\end{exercise}

\begin{exercise}[$\star\star\star$]\label{ex:rw05d}
Using Theorem~\ref{thm:duplication} and Theorem~\ref{thm:betagamma}, show that $B(x,x)=2^{1-2x}B\bigl(\frac12,x\bigr)$ for $x>0$. Check the identity at $x=3$.
\end{exercise}

\section{Hints}

\begin{hint}{ex:halfint}
Use induction on $\Gamma(x+1)=x\Gamma(x)$, or the duplication formula at $x=n+\frac12$. For negative arguments, run the functional equation backwards.
\end{hint}
\begin{hint}{ex:expp}
Put $t=x^p$.
\end{hint}
\begin{hint}{ex:sincos}
Put $t=\sin^2\theta$.
\end{hint}
\begin{hint}{ex:sphere}
Differentiate $r^nV_n$. The maximum is read from the ratio $\frac{S_n}{S_{n-1}}$.
\end{hint}
\begin{hint}{ex:lemniscate}
Substitute $r^4=t$ to get a Beta integral, then use the reflection formula at $\frac14$.
\end{hint}
\begin{hint}{ex:multiplication}
Apply Bohr--Mollerup to $f(x)=\frac{m^{x-1}}C\prod_{k=0}^{m-1}\Gamma\bigl(\frac{x+k}m\bigr)$ with a suitable constant $C$. To find $C$, use
the reflection formula and $\prod_{k=1}^{m-1}\sin\frac{k\pi}m=\frac m{2^{m-1}}$ (Volume~1, Chapter~5, the product of sines).
\end{hint}
\begin{hint}{ex:digamma}
Take the logarithm of Weierstrass's product and differentiate term by term: you need the uniform convergence on $[\frac12,2]$
of the series of derivatives (Lemma~\ref{lem:uniform}).
\end{hint}

\begin{hint}{pb:nair}
(a) Compute $\int_0^1x^{m-1}(1-x)^{n-m}\dd x$ in two ways. (b) $n=2k+1$ and $m=k+1$, then $m=k$ with $n=2k$.
\end{hint}

\begin{hint}{pb:ncycles}
(a) Use the eigenvalues of $c$, and the column orthogonality $\sum_\chi|\chi(c)|^2=|C_G(c)|$. (c) The reciprocal $\frac1{\binom mr}$ is a Beta integral.
\end{hint}

\begin{hint}{pb:dalzell}
(a) Divide the polynomial $x^4(1-x)^4$ by $1+x^2$. (b) $x^4(1-x)^4\equiv-4$ modulo $1+x^2$: evaluate at $x=i$.
\end{hint}

\begin{hint}{pb:sophomore}
Write $x^{-x}=e^{-x\log x}$, expand, and use $\int_0^1x^k(-\log x)^k\dd x=\frac{k!}{(k+1)^{k+1}}$.
\end{hint}

\begin{hint}{ex:rw05a}
$\Gamma(6)=5!$ and $B(3,2)=\Gamma(3)\Gamma(2)/\Gamma(5)$.
\end{hint}

\begin{hint}{ex:rw05b}
$\sin\frac\pi6=\frac12$ and $\sin\frac\pi3=\frac{\sqrt3}2$.
\end{hint}

\begin{hint}{ex:rw05c}
$\Gamma\bigl(\frac12\bigr)=\sqrt\pi$, $\Gamma\bigl(\frac52\bigr)=\frac34\sqrt\pi$ and $\Gamma\bigl(\frac72\bigr)=\frac{15}8\sqrt\pi$. For the ratio of volumes, cancel $\pi^2$.
\end{hint}

\begin{hint}{ex:rw05d}
Duplication gives $\Gamma(2x)$ in terms of $\Gamma(x)\Gamma\bigl(x+\frac12\bigr)$. Then $B\bigl(\frac12,x\bigr)=\sqrt\pi\,\Gamma(x)/\Gamma\bigl(x+\frac12\bigr)$.
\end{hint}

\section{Solutions}

\begin{solution}{ex:halfint}
\idea Each use of the functional equation multiplies by a half-integer. Gather the odd factors into a factorial and the closed form appears.
\writeup By the functional equation,
\[
\begin{aligned}\Gamma\Bigl(n+\frac12\Bigr)&=\frac{2n-1}2\cdot\frac{2n-3}2\cdots\frac12\,\Gamma\Bigl(\frac12\Bigr)\eqbr=\frac{1\cdot3\cdots(2n-1)}{2^n}\sqrt\pi,\end{aligned}
\]
and
$1\cdot3\cdots(2n-1)=\frac{(2n)!}{2\cdot4\cdots2n}=\frac{(2n)!}{2^nn!}$. This gives the formula. Backwards, $\Gamma(-\frac12)=\frac{\Gamma(\frac12)}{-1/2}=-2\sqrt\pi$
and $\Gamma(-\frac32)=\frac{\Gamma(-\frac12)}{-3/2}=\frac43\sqrt\pi$, in accordance with the signs in Figure~\ref{fig:gamma}.
\end{solution}

\begin{solution}{ex:expp}
\idea The substitution $t=x^p$ turns $\int e^{-x^p}x^q\dd x$ into a Gamma integral.
\writeup With $t=x^p$, $x=t^{1/p}$ and $\mathrm dx=\frac1pt^{\frac1p-1}\dd t$, so the integral is $\frac1p\Gamma(\frac1p)=\Gamma(1+\frac1p)$. As $p\to\infty$
it tends to $\Gamma(1)=1$, since $\Gamma$ is continuous. This is obvious from the integrand: $e^{-x^p}$ tends to $1$ on $[0,1[$ and to
$0$ on $]1,\infty[$, so the area tends to that of the unit square. For $p=2$ we recover $\Gamma(\frac32)=\frac{\sqrt\pi}2$.
\end{solution}

\begin{solution}{ex:sincos}
\idea Put $t=\sin^2\theta$. The interval $[0,\frac\pi2]$ goes to $[0,1]$, and $\sin$ and $\cos$ become $t^{1/2}$ and $(1-t)^{1/2}$.
\writeup With $t=\sin^2\theta$, $\mathrm dt=2\sin\theta\cos\theta\dd\theta$, so
\[
\begin{aligned}\int_0^{\pi/2}\sin^a\theta\cos^b\theta\dd\theta&=\frac12\int_0^1t^{\frac{a-1}2}(1-t)^{\frac{b-1}2}\dd t\eqbr=\frac12B\Bigl(\frac{a+1}2,\frac{b+1}2\Bigr).\end{aligned}
\]
With $b=0$, $W_n=\frac12B(\frac{n+1}2,\frac12)=\frac{\sqrt\pi}2\frac{\Gamma(\frac{n+1}2)}{\Gamma(\frac n2+1)}$. For $n=2p$, Exercise~\ref{ex:halfint} gives
$W_{2p}=\frac{(2p)!}{4^p(p!)^2}\frac\pi2$, Wallis's formula. Notice that $W_n=\frac{V_n}{2V_{n-1}}$ (compare with the proof of Theorem~\ref{thm:ballvolume}): Wallis's integrals are
ratios of volumes of balls.
\end{solution}

\begin{solution}{ex:sphere}
\idea A ball is the integral of the areas of the spheres inside it, so the surface area is a derivative of the volume.
\writeup $\frac{\mathrm d}{\mathrm dr}r^nV_n=nr^{n-1}V_n$, which is $nV_n$ at $r=1$. With $\frac n2\Gamma(\frac n2)=\Gamma(\frac n2+1)$ this is
$\frac{2\pi^{n/2}}{\Gamma(n/2)}$. For $n=2$, $\frac{2\pi}{\Gamma(1)}=2\pi$; for $n=3$, $\frac{2\pi^{3/2}}{\sqrt\pi/2}=4\pi$. The ratio
$\frac{S_n}{S_{n-1}}=\sqrt\pi\frac{\Gamma(\frac n2)}{\Gamma(\frac{n+1}2)}$ decreases (by log-convexity, $\frac{\Gamma(y+\frac12)}{\Gamma(y)}$ increases), and it equals
$\frac{33.07}{31.01}>1$ for $n=6$ and $\frac{32.47}{33.07}<1$ for $n=7$. So the largest sphere is $S_6$, in $\R^7$, with area
$\frac{16\pi^3}{15}\approx33.07$.
\end{solution}

\begin{solution}{ex:lemniscate}
\idea A fourth root under a square root is a Beta integral at $\frac14$ and $\frac12$. Reflection trades the leftover $\Gamma(\frac34)$ for $\Gamma(\frac14)$.
\writeup With $r=t^{1/4}$, $\mathrm dr=\frac14t^{-3/4}\dd t$, the integral is $\frac14B(\frac14,\frac12)=\frac{\Gamma(\frac14)\Gamma(\frac12)}{4\Gamma(\frac34)}$. Reflection at
$x=\frac14$ gives $\Gamma(\frac34)=\frac{\pi\sqrt2}{\Gamma(1/4)}$, so the integral is
\[
\frac{\Gamma(\frac14)^2\sqrt\pi}{4\pi\sqrt2}=\frac{\Gamma(\frac14)^2}{4\sqrt{2\pi}}\approx1.31103.
\]
So $\varpi\approx2.62206$. Chapter~\ref{ch:agm} computes $\varpi$ (and so $\Gamma(\frac14)$) by Gauss's arithmetic--geometric mean.
\end{solution}

\begin{solution}{ex:multiplication}
\idea As a function of $mx$, the left side obeys the functional equation of $\Gamma$ up to the factor $m^x$. Bohr--Mollerup names it up to a constant, and reflection fixes that constant.
\writeup Let $C=\prod_{k=1}^{m-1}\Gamma(\frac km)$ and
\[
f(x)=\frac{m^{x-1}}C\prod_{k=0}^{m-1}\Gamma\Bigl(\frac{x+k}m\Bigr).
\]
Then $f>0$, $f(1)=\frac1C\prod_{k=1}^m\Gamma(\frac km)=1$, and
$f$ is log-convex as a product of log-convex functions. Replacing $x$ by $x+1$ shifts the arguments $\frac xm,\dots,\frac{x+m-1}m$ to
$\frac{x+1}m,\dots,\frac{x+m}m$; the product gains the factor $\frac{\Gamma(\frac xm+1)}{\Gamma(\frac xm)}=\frac xm$, and $m^{x-1}$ gains a factor $m$. So
$f(x+1)=xf(x)$, and $f=\Gamma$ by Theorem~\ref{thm:bohrmollerup}.

The constant $C$ comes from pairing $k$ with $m-k$ and using reflection:
\[
C^2=\prod_{k=1}^{m-1}\Gamma\Bigl(\frac km\Bigr)\Gamma\Bigl(1-\frac km\Bigr)=\frac{\pi^{m-1}}{\prod_{k=1}^{m-1}\sin\frac{k\pi}m}.
\]
Recall the proof of $\prod_{k=1}^{m-1}\sin\frac{k\pi}m=\frac m{2^{m-1}}$: with $\omega=e^{2\pi i/m}$, evaluate $\prod_{k=1}^{m-1}(X-\omega^k)=1+X+\cdots+X^{m-1}$ at
$X=1$ and use $|1-e^{i\theta}|=2\sin\frac\theta2$ for $0<\theta<2\pi$. Hence $C^2=\frac{(2\pi)^{m-1}}m$ and $C=\frac{(2\pi)^{(m-1)/2}}{\sqrt m}$, since $C>0$. Now replace $x$
by $mx$ in $\Gamma=f$:
\[
\begin{aligned}\Gamma(mx)&=\frac{m^{mx-1}\sqrt m}{(2\pi)^{(m-1)/2}}\prod_{k=0}^{m-1}\Gamma\Bigl(x+\frac km\Bigr),\end{aligned}
\]
which rearranges to the formula. For $m=2$ it is Legendre's duplication formula.
\end{solution}

\begin{solution}{ex:digamma}
\idea Take the logarithm of Weierstrass's product. The resulting series has a derivative at $1$ that telescopes.
\writeup Taking logarithms in Corollary~\ref{cor:weierstrass},
\[
\begin{gathered}-\ln\Gamma(x)=\ln x+\gamma x+\sum_{k\ge1}u_k(x),\qq u_k(x)=\ln\Bigl(1+\frac xk\Bigr)-\frac xk.\end{gathered}
\]
On $[\frac12,2]$, $u_k'(x)=\frac1{k+x}-\frac1k=-\frac x{k(k+x)}$ satisfies $|u_k'(x)|\le\frac2{k^2}$, so the series of derivatives converges uniformly
(Weierstrass test), and the series itself converges at $x=1$. By Lemma~\ref{lem:uniform} the sum is differentiable with the
termwise derivative, so
\[
\begin{aligned}-\frac{\Gamma'(x)}{\Gamma(x)}&=\frac1x+\gamma+\sum_{k\ge1}\Bigl(\frac1{k+x}-\frac1k\Bigr).\end{aligned}
\]
At $x=1$ the series telescopes to $\sum_k(\frac1{k+1}-\frac1k)=-1$, so $-\Gamma'(1)=1+\gamma-1=\gamma$. By
Proposition~\ref{prop:gammasmooth}, $\Gamma'(1)=\int_0^\infty e^{-t}\ln t\dd t$.
\end{solution}

\begin{solution}{pb:nair}
\idea Expand $\int_0^1x^{m-1}(1-x)^{n-m}\dd x$. The result is a sum of fractions $\frac{\pm c}{m+k}$ whose denominators are at most $n$, so $d_n$ times the integral is an integer. The Beta function evaluates the same integral as $\frac1{m\binom nm}$. Hence $m\binom nm\mid d_n$, and the central binomial coefficient supplies the exponential lower bound.

\writeup (a) Expand $(1-x)^{n-m}$:
\[
I=\int_0^1x^{m-1}(1-x)^{n-m}\dd x=\sum_{k=0}^{n-m}\binom{n-m}k\frac{(-1)^k}{m+k},
\]
and $m+k\le n$, so $d_nI\in\Z$. By the Beta integral (Theorem~\ref{thm:betagamma}), $I=\frac{(m-1)!(n-m)!}{n!}=\frac1{m\binom nm}$. So $m\binom nm\mid d_n$.

(b) Take $n=2k$ and $m=k$: $k\binom{2k}k\mid d_{2k}\mid d_{2k+1}$. With $n=2k+1$, $m=k+1$: $(k+1)\binom{2k+1}{k+1}=(2k+1)\binom{2k}k$ divides
$d_{2k+1}$. As $\gcd(k,2k+1)=1$, $k(2k+1)\binom{2k}k\mid d_{2k+1}$. The largest of the $2k+1$ terms $\binom{2k}j$ of $(1+1)^{2k}$ is at
least their mean: $\binom{2k}k\ge\frac{4^k}{2k+1}$. So $d_{2k+1}\ge k\,4^k$. For $k\ge2$, $k4^k\ge2^{2k+1}$: $d_n\ge2^n$ for odd $n\ge5$. For
even $n=2k+2$, $d_n\ge d_{2k+1}\ge k4^k\ge2^{2k+2}$ as soon as $k\ge4$, that is $n\ge10$; and $d_8=840\ge256$. Together with
$d_7=420\ge128$ and $d_9=2520\ge512$, $d_n\ge2^n$ for all $n\ge7$.

At the end, $d_n=\prod_{p\le n}p^{\lfloor\log n/\log p\rfloor}\le n^{\pi(n)}$, because each $p^{\lfloor\log n/\log p\rfloor}\le n$. So
$\pi(n)\log n\ge\log d_n\ge n\log2$.
\tools This chapter (Beta integral); Volume~2, Chapters~1 and~2 (lcm, Chebyshev bounds); Volume~1, Chapter~1 (binomial coefficients).
\leads Nair (1982). Note that $\log d_n=\psi(n)$, Chebyshev's function, so the prime number theorem is equivalent to
$d_n=e^{n(1+o(1))}$. The same integrals, with more variables, are the start of Beukers's proofs of the irrationality of $\zeta(2)$ and
$\zeta(3)$ (1979), where the bound $d_n\le3^n$ for large $n$ is exactly what is needed.
\end{solution}

\begin{solution}{pb:ncycles}
\idea The class formula writes the count as a sum over characters of $\frac{\chi(c)^3}{\chi(1)}$. Only the exterior powers of $V$ survive, and for those $\chi(c)=\pm1$. The sum that remains is a sum of reciprocals of binomial coefficients, and a Beta integral evaluates it.

\writeup (a) On $\C^n$, the eigenvalues of $c$ are the $n$-th roots of unity $\zeta^j$, $0\le j<n$ (the eigenvectors are
$(\zeta^{jk})_k$), and $V$ is the sum of the eigenspaces with $j\ne0$. So $\chi_r(c)=e_r(\zeta,\dots,\zeta^{n-1})$, the elementary
symmetric function, and
\[
\sum_r\chi_r(c)t^r=\prod_{j=1}^{n-1}(1+t\zeta^j)=\frac{1-(-t)^n}{1+t}=\sum_{r=0}^{n-1}(-t)^r,
\]
using $\prod_{j=0}^{n-1}(1+t\zeta^j)=1-(-t)^n$ (substitute $X=-\frac1t$ in $\prod(X-\zeta^j)=X^n-1$). So $\chi_r(c)=(-1)^r$. The
centraliser of $c$ is the cyclic group it generates, of order $n$, and by column orthogonality (Volume~3, Chapter~11)
$\sum_\chi|\chi(c)|^2=n$. The $n$ characters $\chi_r$ already contribute $n$, so all other irreducible characters vanish at $c$.

(b) Let $\rho$ be irreducible of degree $d$ and let $A$ be a class. The matrix $\sum_{a\in A}\rho(a)$ commutes with $\rho(G)$, so it is the
scalar $\frac{|A|\chi(a)}d$ (compare traces). Then
\[
\#\{(a,b)\in A\times B:ab=g\}=\frac1{|G|}\sum_\chi\chi(1)\sum_{a,b}\chi\bigl(ab\,g^{-1}\bigr),
\]
because $\frac1{|G|}\sum_\chi\chi(1)\chi(h)$ is $1$ if $h=e$ and $0$ otherwise (the regular character). And
$\sum_{a,b}\chi(abg^{-1})=\operatorname{tr}\bigl(\sum_a\rho(a)\sum_b\rho(b)\rho(g)^{-1}\bigr)=\frac{|A|\chi(a)}d\frac{|B|\chi(b)}d\chi(g^{-1})$. Since
$\chi(g^{-1})=\overline{\chi(g)}$, the formula follows.

(c) The Beta integral (Theorem~\ref{thm:betagamma}) gives $\frac1{\binom mr}=(m+1)\int_0^1x^r(1-x)^{m-r}dx$. Sum the geometric progression:
$\sum_{r=0}^m(-x)^r(1-x)^{m-r}=(1-x)^{m+1}-(-x)^{m+1}$, since the ratio is $\frac{-x}{1-x}$ and $(1-x)-(-x)=1$. Integrating,
$\frac1{m+2}+\frac{(-1)^m}{m+2}$; multiply by $m+1$.

(d) Let $A=B$ be the class of $c$, of size $(n-1)!$, and set $g=c$. Only the $\chi_r$ contribute, with $\chi_r(c)^3=(-1)^r$ and
$\chi_r(1)=\binom{n-1}r$:
\[
\frac{(n-1)!^2}{n!}\sum_{r=0}^{n-1}\frac{(-1)^r}{\binom{n-1}r}=\frac{(n-1)!}n\cdot\frac n{n+1}\bigl(1+(-1)^{n-1}\bigr),
\]
which is $\frac{2(n-1)!}{n+1}$ for $n$ odd and $0$ for $n$ even.
\tools Volume~3, Chapter~11 (characters, Schur's lemma, orthogonality); this chapter (Beta integral); Volume~1, Chapter~7 (eigenvalues).
\leads For $n$ even the answer is obvious: an $n$-cycle is odd and a product of two odd permutations is even. For $n=5$ there are $8$
factorisations, for $n=7$, $180$. Stanley (1981) and Zagier (1995) studied the number of cycles of the product of two random
permutations from given classes, a question that is now part of the theory of maps on surfaces.
\end{solution}

\begin{solution}{pb:dalzell}
\idea Divide the numerator by $1+x^2$. A polynomial remains, and its integral is rational. The constant remainder, integrated against $\frac1{1+x^2}$, is a multiple of $\frac\pi4$.

\writeup (a) The division identity is $x^4(1-x)^4=(x^6-4x^5+5x^4-4x^2+4)(1+x^2)-4$. So the integral is
$\frac17-\frac46+1-\frac43+4-4\cdot\frac\pi4=\frac{22}7-\pi$. Since $1\le1+x^2\le2$ on $[0,1]$ and $\int_0^1x^4(1-x)^4\dd x=\frac{4!\,4!}{9!}=\frac1{630}$ (Beta
integral), the integral lies strictly between $\frac1{1260}$ and $\frac1{630}$.

(b) Substitute $x=i$: $x^4(1-x)^4=(1-i)^4=(-2i)^2=-4$, and the same value appears at $x=-i$. So the real polynomial $x^4(1-x)^4+4$ vanishes at $\pm i$ and is
divisible by $1+x^2$. Writing $A=x^4(1-x)^4$, $A^n-(-4)^n=(A+4)(A^{n-1}+\cdots+(-4)^{n-1})$ is divisible by $1+x^2$ too: $A^n=(-4)^n+(1+x^2)Q_n$ with
$Q_n\in\Z[x]$. So the integral equals $r_n+(-4)^n\frac\pi4$ with $r_n=\int_0^1Q_n\in\Q$. It is positive and at most $\int_0^1(x(1-x))^{4n}\le4^{-4n}$.
Multiplying by $4(-4)^{-n}$, $|\pi+4(-4)^{-n}r_n|\le4\cdot4^{-n}\cdot4^{-4n}=4^{1-5n}$, with strict inequality.

\tools{This chapter (Beta integral); Volume~1, Chapter~6 (division of polynomials).}
\leads Dalzell published part (a) in 1944. Stephen Lucas (2009) studied the whole family; the approximations in (b) converge
geometrically, but the denominators of $r_n$ grow too fast for an irrationality proof. Mahler (1953) showed by much finer integrals
that $|\pi-\frac pq|>q^{-42}$ for large $q$; the best exponent known today is $7.10\ldots$ (Zeilberger and Zudilin, 2020).
\end{solution}

\begin{solution}{pb:sophomore}
\idea Write $x^{-x}=e^{-x\log x}$. Each term of the exponential series integrates to a Gamma value, and the $k!$ cancels.

\writeup (a) On $]0,1]$, expand $x^{-x}=\sum_{k\ge0}\frac{(-x\log x)^k}{k!}$. The function $u(x)=-x\log x$, with $u(0)=0$, is continuous on $[0,1]$ and $0\le u\le\frac1e$ there (its
maximum is at $x=\frac1e$). So the $k$-th term is at most $\frac{e^{-k}}{k!}$ in absolute value, the series converges uniformly
on $[0,1]$ by the Weierstrass test, and we may integrate term by term (Lemma~\ref{lem:uniform}). With $x=e^{-u}$,
\[
\int_0^1x^k(-\log x)^k\dd x=\int_0^\infty u^ke^{-(k+1)u}\dd u=\frac{k!}{(k+1)^{k+1}},
\]
so $\int_0^1x^{-x}\dd x=\sum_{k\ge0}\frac1{(k+1)^{k+1}}$. For $x^x=\sum_{k\ge0}\frac{(x\log x)^k}{k!}$ the terms only change sign, the same bound gives uniform convergence, and term-by-term integration gives $\sum_{k\ge0}\frac{(-1)^k}{(k+1)^{k+1}}$.

(b) After $n=N$ the tail is at most $\sum_{n>N}(N+1)^{-n}\le\frac{(N+1)^{-N-1}}{1-1/(N+1)}$. For $N=9$ this is $1.1\cdot10^{-10}$, too large to guarantee ten
decimals; for $N=10$ it is $3.9\cdot10^{-12}$. Ten terms: $1.2912859970\ldots$

\tools{This chapter (Gamma integral); Chapter~\ref{ch:generating} (uniform convergence and term-by-term integration).}
\leads Johann Bernoulli recorded these identities in 1697. The name ``sophomore's dream'' was coined by J.~Borwein, D.~Bailey and
R.~Girgensohn (2004), for an identity that looks like a student's mistake and is true.
\end{solution}

\begin{solution}{ex:rw05a}
\idea The functional equation multiplies by $1,2,3,4,5$ on the way from $\Gamma(1)$ to $\Gamma(6)$. The Beta value is then a ratio of three of those numbers, and the integral is a polynomial.
\writeup $\Gamma(6)=5\cdot4\cdot3\cdot2\cdot1\cdot\Gamma(1)=120$. By Theorem~\ref{thm:betagamma}, $B(3,2)=\frac{\Gamma(3)\Gamma(2)}{\Gamma(5)}=\frac{2\cdot1}{24}=\frac1{12}$. The integral is $\int_0^1(t^2-t^3)\dd t=\frac13-\frac14=\frac1{12}=0.083333333333$.
\end{solution}

\begin{solution}{ex:rw05b}
\idea Reflection at $x$ multiplies $\Gamma(x)$ by $\Gamma(1-x)$ and returns $\pi/\sin(\pi x)$. The two sines here are standard values.
\writeup At $x=\frac16$, $\sin\frac\pi6=\frac12$, so $\Gamma\bigl(\frac16\bigr)\Gamma\bigl(\frac56\bigr)=\pi\big/\frac12=2\pi=6.283185307180$. At $x=\frac13$, $\sin\frac\pi3=\frac{\sqrt3}2$, so the product is $\pi\big/\frac{\sqrt3}2=\frac{2\pi}{\sqrt3}=3.627598728468$.
\end{solution}

\begin{solution}{ex:rw05c}
\idea Even dimensions use $\Gamma(m+1)=m!$. Odd dimensions use $\Gamma\bigl(\frac12\bigr)=\sqrt\pi$ and the functional equation.
\writeup $V_1=\sqrt\pi\big/\Gamma\bigl(\frac32\bigr)=\sqrt\pi\big/\frac{\sqrt\pi}2=2$. $V_2=\pi/\Gamma(2)=\pi$. $V_3=\pi^{3/2}\big/\Gamma\bigl(\frac52\bigr)$ and $\Gamma\bigl(\frac52\bigr)=\frac32\cdot\frac12\sqrt\pi=\frac34\sqrt\pi$, so $V_3=\pi\cdot\frac43=\frac{4\pi}3$. $V_4=\pi^2/\Gamma(3)=\frac{\pi^2}2$. $V_5=\pi^{5/2}\big/\Gamma\bigl(\frac72\bigr)$ and $\Gamma\bigl(\frac72\bigr)=\frac52\cdot\frac32\cdot\frac12\sqrt\pi=\frac{15}8\sqrt\pi$, so $V_5=\pi^2\cdot\frac8{15}=\frac{8\pi^2}{15}$. $V_6=\pi^3/\Gamma(4)=\frac{\pi^3}6$. $V_{10}=\pi^5/\Gamma(6)=\frac{\pi^5}{120}$. $V_{20}=\pi^{10}/\Gamma(11)=\frac{\pi^{10}}{10!}$. Then
\[
\frac{V_5}{V_4}=\frac{8\pi^2/15}{\pi^2/2}=\frac{16}{15},\qquad\frac{V_6}{V_5}=\frac{\pi^3/6}{8\pi^2/15}=\frac{5\pi}{16}.
\]
Since $\pi<\frac{16}5$, one has $\frac{5\pi}{16}<1$, so $V_6<V_5$. The six decimals are $V_1=2$, $V_2=3.141592653590$, $V_3=4.188790204786$, $V_5=5.263789013914$, $V_{10}=2.550164039877$ and $V_{20}=0.025806891390$. The ratio is $V_5/V_{20}=203.968348391$. Figure~\ref{fig:rw05-balls} plots six of these values.
\end{solution}

\begin{solution}{ex:rw05d}
\idea Duplication replaces $\Gamma(2x)$ by a product $\Gamma(x)\Gamma\bigl(x+\frac12\bigr)$. The Beta function at $\bigl(\frac12,x\bigr)$ has the same ratio of Gamma values, up to the constant $\sqrt\pi$.
\writeup Theorem~\ref{thm:duplication} says $\Gamma(2x)=\Gamma(x)\Gamma\bigl(x+\frac12\bigr)\big/\bigl(2^{1-2x}\sqrt\pi\bigr)$. Therefore
\[
B(x,x)=\frac{\Gamma(x)^2}{\Gamma(2x)}=2^{1-2x}\sqrt\pi\,\frac{\Gamma(x)}{\Gamma\bigl(x+\frac12\bigr)}.
\]
Also $B\bigl(\frac12,x\bigr)=\frac{\Gamma\bigl(\frac12\bigr)\Gamma(x)}{\Gamma\bigl(x+\frac12\bigr)}=\sqrt\pi\,\frac{\Gamma(x)}{\Gamma\bigl(x+\frac12\bigr)}$, so $B(x,x)=2^{1-2x}B\bigl(\frac12,x\bigr)$. At $x=3$, $B(3,3)=\frac{\Gamma(3)^2}{\Gamma(6)}=\frac{4}{120}=\frac1{30}$. On the other side, $\Gamma\bigl(\frac72\bigr)=\frac{15}8\sqrt\pi$ and $B\bigl(\frac12,3\bigr)=\frac{\sqrt\pi\cdot2}{15\sqrt\pi/8}=\frac{16}{15}$, so $2^{1-6}\cdot\frac{16}{15}=\frac1{32}\cdot\frac{16}{15}=\frac1{30}=0.033333333333$.
\end{solution}

\section*{Sources and further reading}

\begin{itemize}
\item E.~Artin, \emph{The Gamma Function}, translated by M.~Butler, Holt, Rinehart and Winston, 1964 (German original 1931).
The book is thirty-nine pages. The proofs of Theorems~\ref{thm:betagamma} and~\ref{thm:duplication} follow it.
\item H.~Bohr and J.~Mollerup, \emph{L\ae{}rebog i matematisk Analyse}, vol.~III, Copenhagen, 1922, Theorem~\ref{thm:bohrmollerup} first appeared in that volume.
\item P.~J.~Davis, \emph{Leonhard Euler's integral: a historical profile of the Gamma function}, American Mathematical Monthly
66 (1959), 849--869.
\item G.~Andrews, R.~Askey and R.~Roy, \emph{Special Functions}, Cambridge University Press, 1999, Chapter~1.
\item R.~Remmert, \emph{Wielandt's theorem about the $\Gamma$-function}, American Mathematical Monthly 103 (1996), 214--220. The note is the complex analogue of Bohr--Mollerup, used in Chapter~\ref{ch:complex}.
\item W.~Fulton and J.~Harris, \emph{Representation Theory: A First Course}, Springer, 1991. Proposition~3.12 is the statement quoted in Exercise~\ref{pb:ncycles}.
\item M.~Nair, \emph{On Chebyshev-type inequalities for primes}, American Mathematical Monthly 89 (1982), 126--129, for
Exercise~\ref{pb:nair}.
\item D.~P.~Dalzell, \emph{On 22/7}, Journal of the London Mathematical Society 19 (1944), 133--134; S.~K.~Lucas,
\emph{Approximations to $\pi$ derived from integrals with nonnegative integrands}, American Mathematical Monthly 116 (2009),
166--172, for Exercise~\ref{pb:dalzell}.
\item J.~Borwein, D.~Bailey and R.~Girgensohn, \emph{Experimentation in Mathematics}, A~K~Peters, 2004, for the name
``sophomore's dream'' in Exercise~\ref{pb:sophomore}.
\item C.~Lanczos, \emph{A precision approximation of the gamma function}, SIAM Journal on Numerical Analysis, Series~B, 1
(1964), 86--96. The paper gives a method for numerical evaluation.
\end{itemize}
